\documentclass[11pt]{article}
\usepackage[margin=1in]{geometry}
\usepackage{amsmath,amssymb,amsthm}
\usepackage[hidelinks]{hyperref}
\newtheorem{theorem}{Theorem}
\newtheorem{lemma}{Lemma}
\newcommand{\norm}[1]{\lVert #1\rVert}
\title{The sharp taxicab Erd\H{o}s--Mordell inequality\\
\large MathDB 341671}
\author{Principia Math in Progress}
\date{12 September 2026}
\begin{document}
\maketitle

\paragraph{Task 1. Solve the original conjecture.}
For an interior point of an arbitrary taxicab triangle, prove that the
sum of its distances to the vertices is at least \(3/2\) times the sum
of its distances to the three side lines. This is Conjecture~1 of
Petrovi\'c, Male\v sevi\'c and Banjac~\cite{source}.
We prove a stronger statement.

\begin{theorem}\label{thm:main}
Let \(A_1A_2A_3\) be a nondegenerate triangle in a real normed plane
and let \(P\) be strictly interior. Put
\[
x_i=\norm{A_i-P},\qquad d_i=\operatorname{dist}(P,L_i),
\]
where \(L_i\) is the line containing the side opposite \(A_i\). Then
\begin{equation}\label{eq:main}
\sum_{i=1}^3x_i\ \geq\ \frac32\sum_{i=1}^3d_i.
\end{equation}
The assertion also holds for distances to the finite side segments.
The constant \(3/2\) is optimal in the taxicab plane.
\end{theorem}

\begin{lemma}\label{lem:algebra}
Suppose \(d_i,x_i>0\), \(x_i\geq d_j\) whenever \(i\ne j\), and
\begin{equation}\label{eq:phi}
\Phi(x):=\sum_{i=1}^3\frac{d_i}{d_i+x_i}\leq1.
\end{equation}
Then \eqref{eq:main} holds.
\end{lemma}
\begin{proof}
Relabel the paired indices so that
\[
(d_1,d_2,d_3)=r(1,u,v),\qquad r>0,\quad0<v\leq u\leq1.
\]
The endpoint bounds give \(x\geq b:=r(u,1,1)\) coordinatewise.
If \(2+u\geq\frac32(1+u+v)\), the result follows immediately.
Otherwise set
\[
D=\frac{u+3v-1}{2}>0.
\]
Suppose the conclusion fails and write \(x=b+e\), \(e_i\geq0\).
Its strict failure is precisely
\begin{equation}\label{eq:budget}
\sum_i e_i<rD.
\end{equation}
At the baseline,
\[
\Phi(b)=\frac1{1+u}+\frac{u}{1+u}+\frac{v}{1+v}
=1+\frac{v}{1+v}.
\]
For every \(y\geq b\), each partial derivative obeys
\begin{equation}\label{eq:slope}
\left|\frac{\partial\Phi}{\partial y_i}(y)\right|
=\frac{d_i}{(d_i+y_i)^2}\leq \frac1{r(1+u)^2}.
\end{equation}
For \(i=1,2\) the baseline values are respectively
\(1/[r(1+u)^2]\) and \(u/[r(1+u)^2]\).
For \(i=3\), use
\[
\frac{v}{r(1+v)^2}\leq\frac1{4r}
\leq\frac1{r(1+u)^2}.
\]
Integrating \eqref{eq:slope} along the segment from \(b\) to \(x\)
and then using \eqref{eq:budget} gives
\[
\Phi(x)>1+\frac{v}{1+v}-\frac{D}{(1+u)^2}.
\]
Finally the identity
\begin{align*}
&2v(1+u)^2-(u+3v-1)(1+v)\\
&\hspace{20mm}=1-u+3v(u-v)+2u^2v\geq0
\end{align*}
implies \(D/(1+u)^2\leq v/(1+v)\).
Thus \(\Phi(x)>1\), contradicting \eqref{eq:phi}.
\end{proof}

\begin{proof}[Proof of Theorem~\ref{thm:main}]
Translate \(P\) to \(0\). Its positive barycentric coordinates satisfy
\[
\sum_i\alpha_i=1,\qquad \sum_i\alpha_i A_i=0.
\]
For each \(i\), the point
\[
Q_i=\frac{\sum_{j\ne i}\alpha_j A_j}{1-\alpha_i}
=-\frac{\alpha_i}{1-\alpha_i}A_i
\]
belongs to the opposite side segment. Hence
\[
d_i\leq\norm{Q_i}
=\frac{\alpha_i}{1-\alpha_i}x_i,
\qquad
\frac{d_i}{d_i+x_i}\leq\alpha_i.
\]
Summation proves \eqref{eq:phi}. Moreover, \(A_i\) is an endpoint
of the side opposite \(A_j\) whenever \(i\ne j\), so \(x_i\geq d_j\).
Strict interiority and nondegeneracy give \(d_i,x_i>0\).
Lemma~\ref{lem:algebra} now applies. The same actual points \(Q_i\)
and endpoints establish both constraints for distances to finite
segments, proving that assertion too.

For sharpness take the taxicab norm \(\norm{(s,t)}=|s|+|t|\) and
\[
P=(0,0),\quad A=(1,-\varepsilon),\quad
B=(-1,-\varepsilon),\quad C=(0,1),\qquad\varepsilon>0.
\]
The barycentric coordinates
\[
\left(\frac1{2(1+\varepsilon)},\
\frac1{2(1+\varepsilon)},\
\frac{\varepsilon}{1+\varepsilon}\right)
\]
are positive, so \(P\) is strictly interior. The distance to a line
\(as+bt+c=0\) in this norm is
\(|c|/\max(|a|,|b|)\) at \(P=0\).
Consequently
\[
\sum_i x_i=3+2\varepsilon,\qquad
\sum_i d_i=\varepsilon+\frac2{1+\varepsilon}.
\]
Their difference in \eqref{eq:main} is
\[
3+2\varepsilon-\frac32\left(\varepsilon+\frac2{1+\varepsilon}\right)
=\frac{\varepsilon(7+\varepsilon)}{2(1+\varepsilon)}>0,
\]
and their ratio tends to \(3/2\) as \(\varepsilon\downarrow0\).
No larger universal constant is possible.
\end{proof}

\paragraph{Scope and prior results.}
The source's Theorem~2 covers particular positions relative to the
taxicab axes; its Conjecture~1 on published page~1318 asks for arbitrary
triangles. No rotation is assumed here. The source's definition uses
side lines. The stray extraction suffix in the MathDB display has no
mathematical content and is absent from the typeset inequality.
The current source and public-claim audit is recorded separately.
Earlier norm-dependent estimates and sharpness examples are credited
in that audit; the present proof supplies the universal lower bound.

\paragraph{Task 2. Determine the optimal constant for a fixed norm.}
For a prescribed centrally symmetric unit ball, determine the infimum
of \(\sum_i x_i/\sum_i d_i\) over triangles and interior points.
Identify geometric conditions under which this infimum exceeds \(3/2\).

\paragraph{Task 3. Quantify near equality.}
For taxicab triangles with \(\sum_i x_i=1\), characterize sequences
whose ratios tend to \(3/2\), and derive a quantitative estimate
relating the deficit to their limiting configurations.

\paragraph{Task 4. Find weighted refinements.}
Use the individual constraints
\(d_i/(d_i+x_i)\leq\alpha_i\) to determine sharp inequalities with
prescribed positive barycentric coordinates or vertex weights.

\paragraph{Task 5. Extend the bound to simplices.}
For a simplex in an \(n\)-dimensional normed space, determine the best
universal ratio between the sum of vertex distances and the sum of
distances to opposite facet hyperplanes. Test the analogous
barycentric and endpoint constraints for sufficiency.

\begin{thebibliography}{9}
\bibitem{source}
M. Petrovi\'c, B. Male\v sevi\'c and B. Banjac,
\emph{On the Erd\H{o}s--Mordell inequality for triangles in taxicab geometry},
J. Math. Inequal. \textbf{14} (2020), 1299--1319,
\href{https://doi.org/10.7153/jmi-2020-14-84}{doi:10.7153/jmi-2020-14-84}.
\end{thebibliography}
\end{document}
